%% ============================================================
%%  7.  ANALYSIS
%% ============================================================
\section{Analysis}
\label{sec:analysis}

\subsection{When Does Intervention Help vs.\ Hurt?}
\label{sec:when}

Our results reveal a clear principle: \emph{intervention quality matters
more than intervention quantity}.
Table~\ref{tab:intervention} shows that Exp2a triggers on only 1.5\% of
turns yet still manages to worsen PER (+9.7\%).
Exp3a triggers on 6.2\% of turns and 39.5\% of debates—more than any
other condition—yet produces lower accuracy than Exp3b despite the same
per-turn protection mechanism.

The failure modes differ by condition:

\textbf{Exp2a (loser resistance without quality filter):}
Adversarial incentives cause agents to persist in incorrect minority
positions as well as correct ones.
Because the resistance scalar does not condition on reasoning quality,
wrong agents become equally entrenched, polluting later rounds with
reinforced incorrect hypotheses.
The net effect is increased PER rather than reduced PER.

\textbf{Exp3a (categorical trigger without diagnostic filter):}
Protecting every minority regardless of debate dynamics creates structural
imbalance.
Exp3a raises mean influence asymmetry to $0.235$, compared to $0.098$ for
the baseline, because the intervention \emph{amplifies} minority influence
even when no genuine domination is occurring.
The result is a debate that alternates between majority pressure and
artificially sustained minority dissent, reducing responsiveness (0.011)
and destabilising Nash stability (0.942).

\textbf{Exp3b and Exp4 (selective, quality-gated):}
By conditioning on diagnostic signals (influence asymmetry $\geq 0.65$ or
balance $\leq 0.30$) and, in Exp4, on predicted reasoning quality, these
conditions protect minorities only when the debate dynamics warrant it.
This selectivity keeps influence asymmetry in check (0.141 for Exp3b vs.\
0.235 for Exp3a) while achieving the largest accuracy gains.

The lesson for practitioners is that intervention design requires both a
\emph{when} criterion (diagnostic trigger) and a \emph{who} criterion
(quality filter); either alone is insufficient.

\subsection{Feature Importance of the Learned Predictor}
\label{sec:feature_importance}

\begin{table}[t]
\centering
\caption{Top features of the logistic minority-correctness predictor
(Exp4), sorted by $|\beta|$.
Positive $\beta$ increases probability of correctness; negative $\beta$
decreases it.}
\label{tab:features}
\small
\begin{tabular}{lrl}
\toprule
\textbf{Feature} & $\beta$ & \textbf{Interpretation} \\
\midrule
Minority quality mean     & $+0.245$ & Better reasoning → more likely correct \\
Support delta             & $+0.179$ & Gaining support → more likely correct \\
Num.\ defections          & $-0.121$ & More agents left → less likely correct \\
Was majority before       & $-0.107$ & Defectors to minority → less likely correct \\
Responsiveness            & $-0.098$ & High responsiveness → minority not isolated \\
Balance                   & $+0.091$ & Higher balance → minority survives fairly \\
Rounds remaining          & $-0.088$ & Late rounds → less time to recover \\
Minority confidence mean  & $+0.078$ & Higher self-confidence → slight signal \\
Influence asymmetry       & $+0.077$ & High asymmetry → minority may be suppressed \\
\bottomrule
\end{tabular}
\end{table}

Table~\ref{tab:features} shows that \emph{reasoning quality is the dominant
predictor} ($|\beta|=0.245$), confirming that the LLM judge score carries
genuine signal beyond what process metrics alone provide.
Support trajectory ($\beta=0.179$) indicates that minorities actively gaining
support are more likely correct than those whose support is declining.
Notably, agents who \emph{were} previously in the majority and moved to a
minority position ($\beta=-0.107$) are less likely to be correct—they
likely defected for social rather than epistemic reasons.
High influence asymmetry alone ($\beta=+0.077$) is a weak but positive
predictor: minority agents under a dominant opponent face suppression, and
some fraction of that suppression targets correct views.

Together these features validate the deliberative framework: process
diagnostics that operationalize the Non-domination and Pluralism axioms
carry genuine information about reasoning correctness, not merely about
debate dynamics.

\subsection{Exp4 vs.\ Exp3b: Hard Thresholds vs.\ Learned Predictor}
\label{sec:exp4_vs_exp3b}

\PLACEHOLDER{Fill in this subsection when Exp4 results are available.
Key comparisons to report:
(1) Accuracy and PER of Exp4 vs. Exp3b;
(2) Intervention rate comparison (is Exp4 more selective?);
(3) Are the cases where Exp4 helps / hurts relative to Exp3b interpretable
    from feature importance?
(4) Does scaling tokens by p(correct) help vs. fixed +110 tokens?}
